% GENERATED FILE. Edit canonical lessons, then run npm run generate:books.
% canonical-source-hash: fnv1a64:3d067b42622745a8
% canonical-lessons: TE-C523-karmikudu, TE-C523-nipunudu, TE-C523-salahadaru, TE-C523-sahayakudu, TE-C523-guttedaru

\chapter{People at Work}
\label{ch:te-people-at-work}
\hlchaptermodality{523}

\begin{chapteropening}
\textbf{By the end of this chapter:} \emph{I can name workers, experts, advisers and assistants.}

Complete the last lesson of chapter 523: I can name workers, experts, advisers and assistants.
\end{chapteropening}

\section[\texorpdfstring{\te{కార్మికుడు} (kārmikuḍu)}{కార్మికుడు (kārmikuḍu)}]{\te{కార్మికుడు} (kārmikuḍu) — a worker; a labourer}

\label{lesson:TE-C523-karmikudu}



\begin{quote}
Before the new one: say the Telugu for a promotion, then the Telugu for an ambition.
\end{quote}

\subsection*{You'll want to know: \te{కార్మికుడు}}
\textbf{\te{కార్మికుడు}} — \emph{kārmikuḍu} — \textquotedblleft{}a worker; a labourer\textquotedblright{}.

\begin{quote}
\te{కార్మికులు} \te{ఎండలో} \te{రోడ్డు} \te{వేశారు}.
\end{quote}

\begin{quote}
\emph{kārmikulu eṇḍalō rōḍḍu vēśāru.}
\end{quote}

\begin{quote}
The workers laid the road in the hot sun.
\end{quote}

Plural: \te{కార్మికులు}, \emph{kārmikulu}.

A factory or industrial worker.

\begin{cousinweb}
\textbf{\te{కార్మికుడు}} (\emph{kārmikuḍu}) is built on Sanskrit \emph{karman}, \textquotedblleft{}act, work\textquotedblright{} (CDIAL 2892), with the Telugu masculine ending -\te{డు}: one who works.
\end{cousinweb}

\begin{grammarlens}[title={What you've built}]
One more word for describing what you lived and felt.
\end{grammarlens}

\subsection*{Your turn}
\begin{itemize}
\raggedright
  \item \emph{Say it:} \emph{kārmikuḍu}
  \item \emph{Say it:} \emph{kārmikuḍu}, once more
  \item \emph{Say it:} \emph{kārmikulu eṇḍalō rōḍḍu vēśāru.} — The workers laid the road in the hot sun.
  \item \emph{Recall:} say the Telugu for a promotion, then the Telugu for an ambition, then say \emph{kārmikuḍu} again
\end{itemize}

\begin{tcolorbox}[breakable,colback=teal!4,colframe=teal!35!black,title={Before you move on}]
What does \te{కార్మికుడు} mean? (\textquotedblleft{}A worker; a labourer\textquotedblright{}.) Say it once more.
\end{tcolorbox}
\section[\texorpdfstring{\te{నిపుణుడు} (nipuṇuḍu)}{నిపుణుడు (nipuṇuḍu)}]{\te{నిపుణుడు} (nipuṇuḍu) — an expert; a specialist}

\label{lesson:TE-C523-nipunudu}



\begin{quote}
Before the new one: say the Telugu for an ambition, then the Telugu for a worker.
\end{quote}

\subsection*{You'll want to know: \te{నిపుణుడు}}
\textbf{\te{నిపుణుడు}} — \emph{nipuṇuḍu} — \textquotedblleft{}an expert; a specialist\textquotedblright{}.

\begin{quote}
\te{గుండె} \te{జబ్బుల} \te{నిపుణుడు} \te{నన్ను} \te{పరీక్షించాడు}.
\end{quote}

\begin{quote}
\emph{guṇḍe jabbula nipuṇuḍu nannu parīkṣiñcāḍu.}
\end{quote}

\begin{quote}
A heart specialist examined me.
\end{quote}

Feminine \te{నిపుణురాలు}; \te{నైపుణ్యం}, skill, is the noun.

\begin{cousinweb}
\textbf{\te{నిపుణుడు}} (\emph{nipuṇuḍu}) is borrowed from Sanskrit \emph{nipuṇa}, \textquotedblleft{}clever, adroit, skilful\textquotedblright{} (Monier-Williams), with the Telugu masculine ending -\te{డు}.
\end{cousinweb}

\begin{grammarlens}[title={What you've built}]
One more word for describing what you lived and felt.
\end{grammarlens}

\subsection*{Your turn}
\begin{itemize}
\raggedright
  \item \emph{Say it:} \emph{nipuṇuḍu}
  \item \emph{Say it:} \emph{nipuṇuḍu}, once more
  \item \emph{Say it:} \emph{guṇḍe jabbula nipuṇuḍu nannu parīkṣiñcāḍu.} — A heart specialist examined me.
  \item \emph{Recall:} say the Telugu for an ambition, then the Telugu for a worker, then say \emph{nipuṇuḍu} again
\end{itemize}

\begin{tcolorbox}[breakable,colback=teal!4,colframe=teal!35!black,title={Before you move on}]
What does \te{నిపుణుడు} mean? (\textquotedblleft{}An expert; a specialist\textquotedblright{}.) Say it once more.
\end{tcolorbox}
\section[\texorpdfstring{\te{సలహాదారు} (salahādāru)}{సలహాదారు (salahādāru)}]{\te{సలహాదారు} (salahādāru) — an adviser; a consultant}

\label{lesson:TE-C523-salahadaru}



\begin{quote}
Before the new one: say the Telugu for a worker, then the Telugu for an expert.
\end{quote}

\subsection*{You'll want to know: \te{సలహాదారు}}
\textbf{\te{సలహాదారు}} — \emph{salahādāru} — \textquotedblleft{}an adviser; a consultant\textquotedblright{}.

\begin{quote}
\te{ఆ} \te{కంపెనీకి} \te{ఆయన} \te{సలహాదారు}.
\end{quote}

\begin{quote}
\emph{ā kampenīki āyana salahādāru.}
\end{quote}

\begin{quote}
He is an adviser to that company.
\end{quote}

\te{సలహా}, advice, + -\te{దారు}, holder.

\begin{cousinweb}
\textbf{\te{సలహాదారు}} (\emph{salahādāru}) is built from \te{సలహా}, \textquotedblleft{}advice\textquotedblright{}, which is usually traced through Hindustani to Arabic, and the Persian ending \emph{-dār}, \textquotedblleft{}one who holds\textquotedblright{}.
\end{cousinweb}

\begin{grammarlens}[title={What you've built}]
One more word for describing what you lived and felt.
\end{grammarlens}

\subsection*{Your turn}
\begin{itemize}
\raggedright
  \item \emph{Say it:} \emph{salahādāru}
  \item \emph{Say it:} \emph{salahādāru}, once more
  \item \emph{Say it:} \emph{ā kampenīki āyana salahādāru.} — He is an adviser to that company.
  \item \emph{Recall:} say the Telugu for a worker, then the Telugu for an expert, then say \emph{salahādāru} again
\end{itemize}

\begin{tcolorbox}[breakable,colback=teal!4,colframe=teal!35!black,title={Before you move on}]
What does \te{సలహాదారు} mean? (\textquotedblleft{}An adviser; a consultant\textquotedblright{}.) Say it once more.
\end{tcolorbox}
\section[\texorpdfstring{\te{సహాయకుడు} (sahāyakuḍu)}{సహాయకుడు (sahāyakuḍu)}]{\te{సహాయకుడు} (sahāyakuḍu) — an assistant; a helper}

\label{lesson:TE-C523-sahayakudu}



\begin{quote}
Before the new one: say the Telugu for an expert, then the Telugu for an adviser.
\end{quote}

\subsection*{You'll want to know: \te{సహాయకుడు}}
\textbf{\te{సహాయకుడు}} — \emph{sahāyakuḍu} — \textquotedblleft{}an assistant; a helper\textquotedblright{}.

\begin{quote}
\te{డాక్టరుకు} \te{ఒక} \te{సహాయకుడు} \te{ఉన్నాడు}.
\end{quote}

\begin{quote}
\emph{ḍākṭaruku oka sahāyakuḍu unnāḍu.}
\end{quote}

\begin{quote}
The doctor has an assistant.
\end{quote}

Feminine \te{సహాయకురాలు}; built on \te{సహాయం}, help.

\begin{cousinweb}
\textbf{\te{సహాయకుడు}} (\emph{sahāyakuḍu}) is built on Sanskrit \emph{sahāya}, \textquotedblleft{}companion, helper\textquotedblright{} (CDIAL 13309), with the Telugu masculine ending -\te{డు}.
\end{cousinweb}

\begin{grammarlens}[title={What you've built}]
One more word for describing what you lived and felt.
\end{grammarlens}

\subsection*{Your turn}
\begin{itemize}
\raggedright
  \item \emph{Say it:} \emph{sahāyakuḍu}
  \item \emph{Say it:} \emph{sahāyakuḍu}, once more
  \item \emph{Say it:} \emph{ḍākṭaruku oka sahāyakuḍu unnāḍu.} — The doctor has an assistant.
  \item \emph{Recall:} say the Telugu for an expert, then the Telugu for an adviser, then say \emph{sahāyakuḍu} again
\end{itemize}

\begin{tcolorbox}[breakable,colback=teal!4,colframe=teal!35!black,title={Before you move on}]
What does \te{సహాయకుడు} mean? (\textquotedblleft{}An assistant; a helper\textquotedblright{}.) Say it once more.
\end{tcolorbox}
\section[\texorpdfstring{\te{గుత్తేదారు} (guttēdāru)}{గుత్తేదారు (guttēdāru)}]{\te{గుత్తేదారు} (guttēdāru) — a contractor}

\label{lesson:TE-C523-guttedaru}



\begin{quote}
Before the new one: say the Telugu for an adviser, then the Telugu for an assistant.
\end{quote}

\subsection*{You'll want to know: \te{గుత్తేదారు}}
\textbf{\te{గుత్తేదారు}} — \emph{guttēdāru} — \textquotedblleft{}a contractor\textquotedblright{}.

\begin{quote}
\te{గుత్తేదారు} \te{కొత్త} \te{ఇంటి} \te{పని} \te{మొదలుపెట్టాడు}.
\end{quote}

\begin{quote}
\emph{guttēdāru kotta iṇṭi pani modalupeṭṭāḍu.}
\end{quote}

\begin{quote}
The contractor began work on the new house.
\end{quote}

The one who takes on a job for a fixed price: \te{గుత్త}, a contract.

\begin{cousinweb}
\textbf{\te{గుత్తేదారు}} (\emph{guttēdāru}) is built from \te{గుత్త}, \textquotedblleft{}a contract, a lease\textquotedblright{}, and the ending -\te{దారు}, from Persian \emph{-dār}, \textquotedblleft{}one who holds\textquotedblright{}; the origin of \te{గుత్త} is not settled.
\end{cousinweb}

\begin{grammarlens}[title={What you've built}]
One more word for describing what you lived and felt.
\end{grammarlens}

\subsection*{Your turn}
\begin{itemize}
\raggedright
  \item \emph{Say it:} \emph{guttēdāru}
  \item \emph{Say it:} \emph{guttēdāru}, once more
  \item \emph{Say it:} \emph{guttēdāru kotta iṇṭi pani modalupeṭṭāḍu.} — The contractor began work on the new house.
  \item \emph{Recall:} say the Telugu for an adviser, then the Telugu for an assistant, then say \emph{guttēdāru} again
\end{itemize}

\begin{tcolorbox}[breakable,colback=teal!4,colframe=teal!35!black,title={Before you move on}]
What does \te{గుత్తేదారు} mean? (\textquotedblleft{}A contractor\textquotedblright{}.) Say it once more.
\end{tcolorbox}
